% Clase 20 — Domínguez Palacios Jesús Alejandro (pantalla y RetroArch)
\paragraph{Pantalla}
\textbf{Autor:} Domínguez Palacios Jesús Alejandro

En la pantalla los tipos son formatos de píxel. RetroArch dibuja en un \textit{framebuffer} DRM de
$480 \times 320$ en XRGB8888, un \texttt{uint32\_t} por píxel: el registro de RetroArch reporta
\texttt{stride: 1920}, que es $480 \times 4$ bytes por fila, o sea $480 \times 320 \times 4 = 614\,400$ B por cuadro.
El driver \texttt{panel-mipi-dbi} lo convierte a RGB565 (16 bits por píxel) antes de mandarlo por SPI, así que cada
cuadro que viaja a la pantalla mide $480 \times 320 \times 2 = 307\,200$ B, la mitad.

Con el SPI a 50 MHz (\texttt{config\_txt\_pantalla.txt}, línea 3):
\[
t_{\text{cuadro}} = \frac{307\,200 \times 8\text{ bits}}{50 \times 10^{6}\text{ bit/s}} = 49.15\text{ ms}
\quad\Rightarrow\quad 20.35\text{ FPS máximo}
\]
Los 18.3 FPS medidos en el Avance 2 son el 89.9\,\% de ese máximo: el límite de la pantalla es el bus SPI.

\texttt{volatile} no aparece en esta parte porque no hay código C propio del lado de la Pi: el acceso a los
registros del SPI y del panel lo hace el driver del kernel.
